\section{Joint coefficient laws and nonlinear composition}
\label{jct:section}

The boundary classification in Section~\ref{qfb:section} separates
fixed-order coefficient regimes. We now allow the order and feedback to
move together through the order-three boundary. Exact Lagrange inversion
turns the coefficients into a positive compound-Poisson law; its Fourier
integral produces a common sequence of Gaussian order kernels. This
connects the marginal Lambert coefficients to a uniform moving-order
expansion and to its positive-tail transition.

The joint-coefficient continuation \cite{jct:report} supplies the complete
fixed-logarithmic-order theorem and the positive-tail matching below.
Its two-derivative composition principle also motivates a further result:
Theorem~\ref{jct:thm:matched-jets} isolates the first differing higher
endpoint derivative by comparing two models. This supplies an exact
hierarchy of spatially differentiated correction kernels beyond the
common logarithmic expansion.

Throughout the coefficient discussion, $L=\log n$, $s=3+\eps$ and
$\lambda=\eps L/2$ are real, with $|\lambda|\le\Lambda<\infty$ fixed.
The feedback parameter lies in a fixed compact interval
$[b_-,b_+]\subset(0,\infty)$. Constants may depend on these compact
sets and the fixed truncation order, but not on the varying parameters
or on a lattice location when uniformity in that location is asserted.
The logarithm is principal, with Fourier convention
\begin{equation}\label{jct:eq:fourier-log}
 \operatorname{Log}(-iy)=\log|y|-\frac{\pi i}{2}\operatorname{sgn}(y)
 \qquad(y\in\R\setminus\{0\}).
\end{equation}
We write $\mathfrak g(\xi)=(2\pi)^{-1/2}e^{-\xi^2/2}$ for the normal
density. Square brackets in $\mathcal H_a^{[j]}$ indicate order
derivatives, and $\jctstir{j}{k}$ denotes a Stirling number of the
second kind. Every truncation order is fixed as $n\to\infty$.

\subsection{The exact coefficient law and the resonant scale}
\label{jct:sec:model}

\subsubsection{Quadratic feedback as a compound-Poisson coefficient}

For real $s>2$ and $b>0$, set
\begin{equation}\label{jct:eq:parameters}
 u=\sqrt{b^2+2\zeta(s)},\qquad
 A=\zeta(s-1),\qquad v_c=u-b,\qquad R_b=e^{-v_c}.
\end{equation}
Consider the unique germ
\begin{equation}\label{jct:eq:model}
 uv(r)-\frac{v(r)^2}{2}=\Li_s(re^{v(r)}),\qquad v(0)=0,
 \qquad a_n=[r^n]v(r).
\end{equation}
Define
\begin{equation}\label{jct:eq:f}
 f(z)=u-\sqrt{u^2-2\Li_s(z)}=\sum_{k\ge1}c_kz^k,
\end{equation}
using the square root equal to $u$ at zero. Since
$2\zeta(s)<u^2$, the Taylor expansion of the outer square root converges
absolutely on the whole image of $|z|\le1$. In particular,
\begin{equation}\label{jct:eq:positive-intensity}
 c_k\ge\frac1{u k^s}>0,\qquad
 \sum_{k\ge1}c_k=f(1)=v_c,\qquad f'(1)=\frac Ab.
\end{equation}

\begin{proposition}[Exact coefficient--probability identity]
\label{jct:prop:poisson}
Let $X_n$ be the nonnegative integer-valued compound-Poisson variable
with probability generating function
\begin{equation}\label{jct:eq:pgf}
 \E z^{X_n}=\exp\{n(f(z)-f(1))\}.
\end{equation}
Then
\begin{equation}\label{jct:eq:probability-identity}
 R_b^n a_n=\frac1n\mathbb{P}(X_n=n),\qquad
 \E X_n=\frac{nA}{b}.
\end{equation}
The identity holds for every $b>0$, including $b<A$.
\end{proposition}

\begin{proof}
Put $z=re^{v}$ in \eqref{jct:eq:model}. Solving the quadratic equation
on its germ gives $v=f(z)$, so $z=re^{f(z)}$. Lagrange inversion gives
\[
 a_n=\frac1n[z^{n-1}]f'(z)e^{nf(z)}
     =\frac1n[z^n]e^{nf(z)}.
\]
For the second equality differentiate $e^{nf(z)}$ and compare its
$z^{n-1}$ coefficient. Positivity and summability in
\eqref{jct:eq:positive-intensity} make
$\exp\{n\sum c_k(z^k-1)\}$ a probability generating function:
take a Poisson number of jumps of total intensity $nf(1)$, with
jump probabilities $c_k/f(1)$. Multiplication by $e^{-nf(1)}$
proves the coefficient identity. Differentiation at $z=1$ proves
the mean formula.
\end{proof}

The endpoint reference $R_b$ is the actual Taylor radius at and above
$b=A$. Below that threshold the branch first encounters an interior
fold and the radius is larger than $R_b$ \cite{qfb:reports}.
No radius claim is needed for \eqref{jct:eq:probability-identity}.

\subsubsection{Pairing the colliding Gamma and zeta terms}

Write $s=3+\eps$ and define
\begin{align}
 E_2(\eps)&=\frac{\Gamma(1-\eps)}{(1+\eps)(1+\eps/2)},
 & Z(\eps)&=\zeta(1+\eps)-\eps^{-1},\label{jct:eq:E-Z}\\
 K_\eps(\ell)&=Z(\eps)+
      \frac{1-E_2(\eps)e^{-\eps\ell}}{\eps}.
      \label{jct:eq:K}
\end{align}
Both expressions are interpreted by removal of the singularity at
$\eps=0$. The elementary values are
\begin{equation}\label{jct:eq:resonant-values}
 E_2(0)=1,\quad E_2'(0)=\gamma-\frac32,\quad
 Z(0)=\gamma,\quad K_0(\ell)=\ell+\frac32.
\end{equation}
Gamma recurrence gives
$2\Gamma(-2-\eps)=-E_2(\eps)/\eps$. Pairing that term with the
zeta pole in the convergent Jonqui\`ere expansion yields the following
uniform local statement.

\begin{lemma}[The centered local exponent]
\label{jct:lem:local}
There are $r>0$ and $C<\infty$, uniform for $b\in[b_-,b_+]$ and
real $|\eps|\le1/4$, such that, for
$0<|t|<r$ and $|\arg t|\le\pi/2$,
\begin{equation}\label{jct:eq:local-exponent}
 f(e^{-t})-f(1)+\frac Ab t
 =\frac{t^2}{2b}\left\{K_\eps(-\operatorname{Log} t)+\nu\right\}
   +\mathcal E(t),\qquad \nu=\frac{A^2}{b^2},
\end{equation}
where $|\mathcal E(t)|\le C|t|^{5/2}$.
\end{lemma}

\begin{proof}
The classical normally convergent expansion \cite[25.12.12]{rw:DLMFPolylog}
gives, after pairing,
\begin{equation}\label{jct:eq:deficit-pair}
 D_s(t):=\zeta(s)-\Li_s(e^{-t})
 =At-\frac{t^2}{2}K_\eps(-\operatorname{Log} t)+t^3 J_\eps(t),
\end{equation}
where $J_\eps$ is jointly holomorphic for $t$ near zero and $\eps$
near the indicated compact interval. Normal convergence can also be
checked by the zeta functional equation: the coefficients of the
analytic tail have geometric majorants on each smaller disk
$|t|<r<2\pi$. All poles near $\eps=0$ occur in the displayed paired term.

Using
\[
 \frac{1-e^{-\eps\ell}}{\eps}
  =\ell\int_0^1e^{-\eps\ell v}\,dv
\]
and the bounded regular part of $(E_2(\eps)-1)/\eps$ gives
\[
 |K_\eps(-\operatorname{Log} t)|
 \le C|t|^{-1/4}(1+|\log|t||)\le C'|t|^{-1/2}.
\]
Thus $D_s(t)=At+O(|t|^{3/2})=O(|t|)$. Uniformly in $b$,
\[
 b-\sqrt{b^2+2D}
 =-\frac Db+\frac{D^2}{2b^3}+O(|D|^3).
\]
Inserting \eqref{jct:eq:deficit-pair} leaves the displayed terms of
\eqref{jct:eq:local-exponent}; the difference between $D^2$ and
$A^2t^2$ is $O(|t|^{5/2})$. This proves the bound, including the
imaginary boundary rays by continuation and continuity.
\end{proof}

The nonlinear feedback contributes precisely $\nu$ to the quadratic
regular part. At criticality $b=A$, this is $1$, changing the
resonant constant from $3/2$ to $5/2$.

\subsubsection{A scalar equation absorbing the slow logarithms}

Define $B=B(n,\eps,b)>1$ by
\begin{equation}\label{jct:eq:scale}
 \boxed{\displaystyle\quad
 V:=\frac{bB^2}{n}=K_\eps(\log B)+\nu.
 \quad}
\end{equation}
This defines a scale from the paired quadratic exponent; it does not
assume that the underlying distribution has a finite variance. Put
\begin{equation}\label{jct:eq:q-xi}
 q=\frac{E_2(\eps)B^{-\eps}}{V},\qquad
 \xi_m=\frac{m-nA/b}{B},\qquad \xi=\xi_n.
\end{equation}

\begin{theorem}[Existence, uniqueness, and size of the scale]
\label{jct:thm:scale}
Uniformly for $|\lambda|\le\Lambda$ and $b\in[b_-,b_+]$,
\eqref{jct:eq:scale} has exactly one solution $B>1$ for all sufficiently
large $n$. Let
\begin{equation}\label{jct:eq:leading-scale}
 \varphi(\lambda)=\frac{1-e^{-\lambda}}{\lambda},\quad \varphi(0)=1,
 \qquad B_0^2=\frac{nL}{2b}\varphi(\lambda).
\end{equation}
Then
\begin{align}
 \frac B{B_0}
 &=1+\frac{C_\lambda}{L\varphi(\lambda)}
       +O\left(\frac{(\log L)^2}{L^2}\right),\label{jct:eq:scale-first}\\
 C_\lambda
 &=\gamma+\nu+e^{-\lambda}
   \left\{\frac12\log L+\frac12\log\frac{\varphi(\lambda)}{2b}
                   -\gamma+\frac32\right\},\label{jct:eq:C-lambda}\\
 q&=\frac{2e^{-\lambda}}{L\varphi(\lambda)}
          +O\left(\frac{\log L}{L^2}\right).
          \label{jct:eq:q-equivalent}
\end{align}
In particular $q>0$ and $q\asymp L^{-1}$ on the specified parameter sets.
\end{theorem}

\begin{proof}
Set $H=\log B$ and
$F(H)=be^{2H}/n-K_\eps(H)-\nu$. For small real $\eps$,
$K_\eps(0)+\nu$ is uniformly positive by
\eqref{jct:eq:resonant-values}. Thus $F(0)<0$ for large $n$.
Moreover
\[
 F'(H)=e^{-\eps H}
 \left\{\frac{2b}{n}e^{(2+\eps)H}-E_2(\eps)\right\}.
\]
The braced expression increases strictly, because $2+\eps>0$.
It follows that $F$ first decreases and then increases, or increases
throughout. In both cases $F(H)\to+\infty$ as $H\to\infty$:
the growth exponent $2$ dominates the possible growth exponent
$-\eps$ of $K_\eps$. Hence there is exactly one positive zero.

Write $H_0=\log B_0$
$=\tfrac12\{L+\log L+\log(\varphi(\lambda)/(2b))\}$.
The representation \eqref{jct:eq:K}, with
$\eps=2\lambda/L$, gives uniformly, including $\lambda=0$,
\[
 K_\eps(H_0)+\nu
 =\frac L2\varphi(\lambda)+C_\lambda
                   +O((\log L)^2/L).
\]
For example the first term follows from
$(1-e^{-\lambda})/\eps=L\varphi(\lambda)/2$; the next term uses
$E_2'(0)=\gamma-3/2$. First bracket the root by $B=cB_0$
for fixed constants $c<1$ and $c>1$; the leading terms have opposite
signs. Then the derivative at the bracketed root is of order $L$.
Taylor's theorem in $H-H_0$ improves the bracket to
$H-H_0=C_\lambda/(L\varphi(\lambda))+O((\log L)^2/L^2)$.
Exponentiation gives \eqref{jct:eq:scale-first}. Finally
$E_2(\eps)B^{-\eps}=e^{-\lambda}(1+O(\log L/L))$
and $V=L\varphi(\lambda)/2+O(\log L)$, proving the $q$ formula.
\end{proof}

\subsection{Gaussian order derivatives and finite correction formulas}
\label{jct:sec:kernels}

For $\Real a>-1$ and real $\xi$, define
\begin{equation}\label{jct:eq:H-definition}
 \mathcal{H}_a(\xi)=\frac1{2\pi}\int_{\R}
     e^{-i\xi y-y^2/2}(-iy)^a\,dy,
 \qquad \mathcal{H}_a^{[j]}(\xi)=\partial_a^j\mathcal{H}_a(\xi).
\end{equation}
The logarithm convention is \eqref{jct:eq:fourier-log}. All fixed order
derivatives are legitimate by domination at zero and infinity. For
real $a$ and $\xi$ the two halves of the integral are conjugates.

\begin{proposition}[Parabolic cylinder and Gamma forms]
\label{jct:prop:H}
In the standard parabolic-cylinder normalization,
\begin{align}
 \mathcal{H}_a(\xi)&=\frac{e^{-\xi^2/4}}{\sqrt{2\pi}}D_a(-\xi),
       \label{jct:eq:H-cylinder}\\
 \mathcal{H}_a(0)&=\frac{2^{(a-1)/2}}{\Gamma((1-a)/2)},
 &\partial_\xi\mathcal{H}_a&=\mathcal{H}_{a+1}.
       \label{jct:eq:H-zero}
\end{align}
In particular $\mathcal{H}_0=\mathfrak{g}$ and
$\mathcal{H}_{2k}=\mathfrak{g}^{(2k)}$. Each function
$\mathcal{H}_a^{[j]}$ used below is an ordinary convergent integral, even
when a later asymptotic series containing it does not converge.
\end{proposition}

\begin{proof}
The derivative identity follows immediately under the integral.
Splitting into $y>0$ and $y<0$, then using the Gamma integral,
gives
\[
 \mathcal{H}_a(0)=\frac{2^{(a-1)/2}}{\pi}
       \Gamma((a+1)/2)\cos(\pi a/2).
\]
The reflection formula gives \eqref{jct:eq:H-zero}. Integration by
parts in $y$, with vanishing boundary term
$y(-iy)^ae^{-y^2/2-i\xi y}$, gives
\[
 \mathcal{H}_a''+\xi\mathcal{H}_a'+(a+1)\mathcal{H}_a=0.
\]
The function on the right of \eqref{jct:eq:H-cylinder} satisfies the same
equation by the parabolic-cylinder differential equation.
The values of $D_a(0)$ and $D_a'(0)$
\cite[12.2.4--12.2.7]{jct:cylinder} give the same value and first
derivative at $\xi=0$, namely
$2^{(a-1)/2}/\Gamma((1-a)/2)$ and
$2^{a/2}/\Gamma(-a/2)$. Uniqueness for the ordinary differential
equation proves the identity.
\end{proof}

\subsubsection{The correction polynomials in two small parameters}

Define the finite functions
\begin{align}
 \mathcal C_0(\xi;q,\eps)&=\mathfrak{g}(\xi),\notag\\
 \mathcal C_j(\xi;q,\eps)
 &=\frac1{j!}\sum_{k=1}^j
     \jctstir{j}{k}\left(-\frac q2\right)^k\eps^{j-k}
           \mathcal{H}_{2k}^{[j]}(\xi),\qquad j\ge1.
           \label{jct:eq:correction}
\end{align}
They are homogeneous of degree $j$ in $(q,\eps)$. The first three are
\begin{align}
 \mathcal C_1&=-\frac q2\mathcal{H}_2^{[1]},\label{jct:eq:C1}\\
 \mathcal C_2&=-\frac{q\eps}{4}\mathcal{H}_2^{[2]}
                  +\frac{q^2}{8}\mathcal{H}_4^{[2]},\label{jct:eq:C2}\\
 \mathcal C_3&=-\frac{q\eps^2}{12}\mathcal{H}_2^{[3]}
             +\frac{q^2\eps}{8}\mathcal{H}_4^{[3]}
             -\frac{q^3}{48}\mathcal{H}_6^{[3]}.
             \label{jct:eq:C3}
\end{align}
The spatial argument $\xi$ is suppressed on the right. These expressions
involve order derivatives of the Gaussian kernel; in general they are
not a Gaussian times an ordinary polynomial in $\xi$.

There is also a finite divided-difference form. For a holomorphic
function of $a$, put
\begin{equation}\label{jct:eq:divided-difference}
 \mathfrak D_\eps^k F(a)
 =\frac1{\eps^k}\sum_{h=0}^k(-1)^{k-h}\binom{k}{h}F(a+h\eps),
 \qquad \mathfrak D_0^k F=F^{(k)}.
\end{equation}
The apparent poles are removable. The identity
\begin{equation}\label{jct:eq:stirling-generating}
 \left(\frac{e^{\eps z}-1}{\eps}\right)^k
 =k!\sum_{j\ge k}\jctstir{j}{k}\frac{\eps^{j-k}z^j}{j!}
\end{equation}
explains the appearance of the Stirling numbers. We will prove that
both the $q$-ordered finite expression
\begin{equation}\label{jct:eq:finite-difference-approx}
 \sum_{k=0}^N\frac{(-q/2)^k}{k!}
       \left.\mathfrak D_\eps^k\mathcal{H}_a(\xi)\right|_{a=2k}
\end{equation}
and the total-degree expression
$\sum_{j=0}^N\mathcal C_j(\xi;q,\eps)$ approximate the same lattice
law with error $O(L^{-N-1})$. They differ from one another only at
the omitted orders. Formula \eqref{jct:eq:finite-difference-approx}
can be useful for computation away from $\eps=0$; near zero it is
better evaluated using derivatives or the removable integral formula.

\subsubsection{Finite Gamma moments at the center}

For integers $j,k\ge1$ set
\begin{equation}\label{jct:eq:M-jk}
 M_j(k)=\Real\left[
   (\partial_z+\log2+i\pi)^j\Gamma(z)
                      \right]_{z=k+1/2}.
\end{equation}
This is also the convergent moment
\begin{equation}\label{jct:eq:M-integral}
 M_j(k)=\Real\int_0^\infty e^{-t}t^{k-1/2}
                          (\log t+\log2+i\pi)^j\,dt.
\end{equation}

\begin{corollary}[Central coefficient calculus]
\label{jct:cor:central-corrections}
For every $j\ge1$,
\begin{equation}\label{jct:eq:central-C}
 \frac{\mathcal C_j(0;q,\eps)}{\mathfrak{g}(0)}
 =\frac1{2^j j!\sqrt\pi}
     \sum_{k=1}^j\jctstir{j}{k}q^k\eps^{j-k}M_j(k).
\end{equation}
In particular
\begin{equation}\label{jct:eq:central-first}
 \frac{\mathcal C_1(0;q,\eps)}{\mathfrak{g}(0)}
 =\frac q4(2-\gamma-\log2).
\end{equation}
\end{corollary}

\begin{proof}
Differentiate the Gamma--cosine expression in the proof of
Proposition~\ref{jct:prop:H}. At $a=2k$ it gives
\[
 \mathcal{H}_{2k}^{[j]}(0)=
       \frac{(-1)^k2^{k-j-1/2}}{\pi}M_j(k).
\]
Substitution in \eqref{jct:eq:correction} cancels the two factors
$(-1)^k$ and proves \eqref{jct:eq:central-C}. For $j=k=1$, use
$\Gamma(3/2)=\sqrt\pi/2$ and
$\psi(3/2)=2-\gamma-2\log2$.
\end{proof}

These formulas can be converted into polygamma values, harmonic numbers,
zeta values, and the Euler constant by repeated differentiation of
$\Gamma$ and the half-integer polygamma recurrences. The moment form
keeps all signs and factors of two in a single finite expression.

\subsection{The uniform coefficient transseries}
\label{jct:sec:main-proof}

We now prove the main result. The definition of $B$ in
\eqref{jct:eq:scale} incorporates all logarithmic changes of the quadratic
coefficient, including the slowly varying feedback contribution.

\begin{theorem}[Every fixed logarithmic order, uniformly on the lattice]
\label{jct:thm:all-orders}
Fix $\Lambda<\infty$, $0<b_-<b_+<\infty$, and an integer $N\ge0$.
For $L=\log n$, $|\eps L/2|\le\Lambda$, and $b\in[b_-,b_+]$,
let $X_n$, $B$, $q$, and $\xi_m$ be defined by
\eqref{jct:eq:pgf}, \eqref{jct:eq:scale}, and \eqref{jct:eq:q-xi}. Then
\begin{equation}\label{jct:eq:uniform-lattice}
 \sup_{m\in\mathbb Z}\left|
 B\mathbb{P}(X_n=m)-\sum_{j=0}^N\mathcal C_j(\xi_m;q,\eps)
 \right|\le C_N L^{-N-1}
\end{equation}
for all sufficiently large $n$, uniformly in the stated parameters.
Consequently the feedback coefficients satisfy
\begin{equation}\label{jct:eq:all-orders-coefficient}
 \boxed{\displaystyle
 nB R_b^n a_n
 =\mathfrak{g}(\xi)+
  \sum_{j=1}^N\frac1{j!}\sum_{k=1}^j
    \jctstir{j}{k}\left(-\frac q2\right)^k\eps^{j-k}
           \mathcal{H}_{2k}^{[j]}(\xi)
       +O(L^{-N-1}).}
\end{equation}
The finite divided-difference expression
\eqref{jct:eq:finite-difference-approx} has the same error bound.
\end{theorem}

The theorem resolves the joint coefficient target Q1 in
\cite{qfb:reports}. It is an absolute estimate. For bounded $\xi$
it gives relative Gaussian asymptotics because $\mathfrak{g}$ has a
positive minimum on each compact real interval. Relative conclusions
for unbounded $\xi$ will be proved separately.

\subsubsection{Uniform control of the whole Fourier contour}

\begin{lemma}[Uniform Fourier domination]
\label{jct:lem:fourier-tail}
Let
\begin{equation}\label{jct:eq:chi}
 \chi_n(y)=\E\exp\left(iy\frac{X_n-nA/b}{B}\right).
\end{equation}
There are positive constants $c_1,c_2$, uniform in the moving-order
window and in $b$, such that
\begin{align}
 |\chi_n(y)|&\le e^{-c_1y^2}
       &&(|y|\le Bn^{-1/4}),\label{jct:eq:fourier-small}\\
 |\chi_n(y)|&\le e^{-c_2\sqrt n}
       &&(Bn^{-1/4}\le |y|\le\pi B).
       \label{jct:eq:fourier-large}
\end{align}
\end{lemma}

\begin{proof}
By positivity of the intensities,
\[
 -\log|\chi_n(B\theta)|
  =n\sum_{k\ge1}c_k(1-\cos k\theta).
\]
For $|\theta|\le n^{-1/4}$ retain $1\le k\le\lfloor n^{1/4}\rfloor$.
Then $|k\theta|\le1$, so $1-\cos k\theta\ge c k^2\theta^2$.
The parameters keep $u$ bounded above. Moreover
$k^{-\eps}\ge e^{-\Lambda/2}$ on this range, since
$|\eps|\log k\le\Lambda/2$. Therefore
\[
 -\log|\chi_n(B\theta)|
 \ge c n\theta^2\sum_{k\le n^{1/4}}k^{-1-\eps}
 \ge c' nL\theta^2.
\]
Since $B^2\asymp nL$, this proves \eqref{jct:eq:fourier-small}.
For $n^{-1/4}\le|\theta|\le\pi$, retain only $k=1$.
Here $c_1=1/u$ in the intensity notation and
$1-\cos\theta\ge2\theta^2/\pi^2$, giving
$-\log|\chi_n(B\theta)|\ge c n\theta^2\ge c\sqrt n$.
The centering phase has modulus one, proving the second assertion.
\end{proof}

This argument is insensitive to the cancellation of the Gamma and
zeta poles. It uses the exact nonnegative coefficients and supplies
the uniformity that a fixed-order local-limit theorem alone would
not provide.

\subsubsection{The exact central reduction}

For $z\in\C$, define the entire divided exponential
\begin{equation}\label{jct:eq:h-epsilon}
 h_\eps(z)=\frac{e^{\eps z}-1}{\eps},\qquad h_0(z)=z.
\end{equation}
Take $t=-iy/B$ in Lemma~\ref{jct:lem:local}. The scale equation removes the
constant quadratic term and gives
\begin{equation}\label{jct:eq:core-characteristic}
 \chi_n(y)=\exp\left\{
     -\frac{y^2}{2}+\frac q2y^2h_\eps(\operatorname{Log}(-iy))+R_n(y)
                       \right\},
 \quad |R_n(y)|\le CnB^{-5/2}|y|^{5/2},
\end{equation}
whenever $0<|y|<rB$. The limit at zero is understood continuously.
The sign in this formula follows from
$(-iy)^2=-y^2$ and
\[
 K_\eps(\log B-\operatorname{Log}(-iy))-K_\eps(\log B)
   =-E_2(\eps)B^{-\eps}h_\eps(\operatorname{Log}(-iy)).
\]
Since $nB^{-5/2}=O(n^{-1/4}L^{-5/4})$, the local nonlinear
remainder is smaller than every fixed inverse power of $L$.

\subsubsection{A finite expansion with an integrable remainder}

\begin{proof}[Proof of Theorem~\ref{jct:thm:all-orders}]
Choose
$T_L=\sqrt{D_N\log L}$ with a sufficiently large constant $D_N$.
By Lemma~\ref{jct:lem:fourier-tail}, the part of the scaled Fourier integral
with $T_L<|y|\le\pi B$ has absolute size
$O(L^{-N-1})$: use the Gaussian bound up to $Bn^{-1/4}$ and
the bound $2\pi B e^{-c_2\sqrt n}$ beyond it.

On $|y|\le T_L$, put $z=\operatorname{Log}(-iy)$ and
$w=(q/2)y^2h_\eps(z)$. The inequality
\begin{equation}\label{jct:eq:h-bound}
 |h_\eps(z)|\le |z|e^{|\eps z|}
\end{equation}
follows from $h_\eps(z)=z\int_0^1e^{v\eps z}\,dv$.
It proves $\sup_{|y|\le T_L}|w|=o(1)$. Near $y=0$ one uses
$y^2|\log|y||\,|y|^{-|\eps|}\to0$; no uniform bound on
$|\eps\log|y||$ is assumed. For $1\le|y|\le T_L$ the bound is
$O(L^{-1}\log L\log\log L)$.
Equation \eqref{jct:eq:core-characteristic} and the estimate on $R_n$
therefore allow us to replace $\chi_n$ by
$e^{-y^2/2}e^w$ in the central integral with error
$O(nB^{-5/2})$.

Expand the exponential only to the finite order $N$:
\[
 e^w=\sum_{k=0}^N\frac{w^k}{k!}+O_N(|w|^{N+1}).
\]
The error is integrable against $e^{-y^2/2}$ with bound
$C_N q^{N+1}$. Indeed \eqref{jct:eq:h-bound} bounds the integrand near
zero by a constant times
$|y|^{2N+2-(N+1)|\eps|}(1+|\log|y||)^{N+1}$;
at infinity a Gaussian dominates the polynomial and logarithmic
factors. For the fixed $N$, take $n$ large enough that
$(N+1)|\eps|<1/2$. All these integrals then have a common majorant.

For each $k\le N$, use
\begin{equation}\label{jct:eq:cube-integral}
 h_\eps(z)^k=z^k\int_{[0,1]^k}
      e^{\eps z(v_1+\cdots+v_k)}\,dv_1\cdots dv_k.
\end{equation}
Taylor expansion in $\eps$ through degree $N-k$ gives
\[
 h_\eps(z)^k
 =k!\sum_{j=k}^N\jctstir{j}{k}\frac{\eps^{j-k}z^j}{j!}
    +\mathcal R_{k,N}(z),
\]
where
\[
 |\mathcal R_{k,N}(z)|\le
 C_N|\eps|^{N+1-k}|z|^{N+1}e^{k|\eps z|}.
\]
After multiplication by $q^k|y|^{2k}e^{-y^2/2}$, its integral is
$O(q^k|\eps|^{N+1-k})=O(L^{-N-1})$, by the same domination.
Thus the total-degree polynomial in $q,\eps$ has an $L^1(dy)$
error of the asserted order. Its finite terms may be extended from
$|y|\le T_L$ to the full real axis: their Gaussian tails are
$O(L^{-N-1})$ after enlarging $D_N$ if needed.

The resulting integrals are
\[
 \frac1{2\pi}\int_\R e^{-i\xi_m y-y^2/2}
       (-iy)^{2k}(\operatorname{Log}(-iy))^j\,dy
       =\mathcal{H}_{2k}^{[j]}(\xi_m).
\]
Here the factor $y^{2k}$ in $w^k$ contributes $(-1)^k$, producing
$(-q/2)^k$. Lattice Fourier inversion says
\[
 B\mathbb{P}(X_n=m)
 =\frac1{2\pi}\int_{-\pi B}^{\pi B}
                 e^{-i\xi_m y}\chi_n(y)\,dy.
\]
Every error bound above is an $L^1$ bound before insertion of the
phase. Since $|e^{-i\xi_m y}|=1$, it is uniform in all integers $m$.
This proves \eqref{jct:eq:uniform-lattice}. The exact coefficient identity
gives \eqref{jct:eq:all-orders-coefficient}.

If one retains $h_\eps(z)^k$ instead of Taylor-expanding in
$\eps$, apply finite Fourier inversion to
\[
 (-iy)^{2k}h_\eps(\operatorname{Log}(-iy))^k.
\]
This gives precisely \eqref{jct:eq:finite-difference-approx}. The same
$q^{N+1}$ estimate proves its assertion.
\end{proof}

\begin{remark}[The full formal exponential is not the Fourier law]
The proof integrates a finite expansion in the central region and
uses the exact characteristic function on its complement. It does not
replace the full Fourier law by
$\exp\{-y^2/2+(q/2)y^2h_\eps(\operatorname{Log}(-iy))\}$ on the whole real axis.
For fixed positive $q,\eps$ the latter expression can grow at large
$|y|$. Nor does the theorem assert convergence of the infinite sum
over $j$. The finite integrals and all fixed-order remainders are the
claimed objects.
\end{remark}

\subsubsection{The elementary normalization requested in Q1}

Put $\xi_0=n(1-A/b)/B_0$, using \eqref{jct:eq:leading-scale}.
On compact $\xi_0$ sets, Theorem~\ref{jct:thm:scale} and Theorem~\ref{jct:thm:all-orders} give
\begin{equation}\label{jct:eq:elementary-first}
\begin{split}
 nB_0R_b^n a_n={}&\mathfrak{g}(\xi_0)
   +\frac{C_\lambda\mathcal{H}_2(\xi_0)
               -e^{-\lambda}\mathcal{H}_2^{[1]}(\xi_0)}{L\varphi(\lambda)}\\
 &+O\left(\frac{(\log L)^2}{L^2}\right).
\end{split}
\end{equation}
To verify the scale-change term, write
$B/B_0=1+\delta$; then
$(B_0/B)\mathfrak{g}(\xi_0B_0/B)
=\mathfrak{g}(\xi_0)-\delta(\mathfrak{g}+\xi_0\mathfrak{g}')(\xi_0)+O(\delta^2)$,
and $-(\mathfrak{g}+\xi\mathfrak{g}')=\mathcal{H}_2$.

In particular, if $\xi_0$ remains bounded, the leading statement is
\begin{equation}\label{jct:eq:Q1-solved}
 n\sqrt{\frac{n\log n}{2b}\varphi(\lambda)}\,R_b^n a_n
 =\mathfrak{g}\left(
 \frac{n(1-\zeta(2+\eps)/b)}{
               \sqrt{n\log n\,\varphi(\lambda)/(2b)}}\right)+o(1).
\end{equation}
At the critical choice $b=A=\zeta(2+\eps)$,
\begin{equation}\label{jct:eq:critical-joint-leading}
 a_n\sim R_b^{-n}
   \sqrt{\frac{A}{\pi\varphi(\lambda)}}\,
                   n^{-3/2}(\log n)^{-1/2},
\end{equation}
uniformly over the bounded moving-order window. This is a joint
transition law; letting $\lambda$ tend to infinity is outside its
stated uniform parameter set.

\subsection{Lambert resummation and nonlinear universality}
\label{jct:sec:universality}

\subsubsection{A convergent chart for the scale itself}

The coefficient expansion is asymptotic. The scalar scale $B$ has a
stronger analytic structure that explains how it absorbs the
$\log\log n$ corrections.

\begin{proposition}[Convergent two-variable scale chart]
\label{jct:prop:scale-chart}
Let $h=L^{-1}$, $w=L^{-1}\log L$, and
$c=\log(\varphi(\lambda)/(2b))$. There is a holomorphic germ
$Y(h,w;\lambda,b)$, with $Y(0,0;\lambda,b)=1$, uniform on compact
real $(\lambda,b)$ sets as above, such that
\begin{equation}\label{jct:eq:Y-scale}
 B^2=B_0^2Y(L^{-1},L^{-1}\log L;\lambda,b).
\end{equation}
It is determined by the removable holomorphic equation
\begin{equation}\label{jct:eq:Y-equation}
\begin{split}
 Y={}&\frac{1-e^{-\lambda}E_2(2\lambda h)
                e^{-\lambda w-\lambda h(c+\log Y)}}{
                \lambda\varphi(\lambda)}\\
 &+\frac{2h}{\varphi(\lambda)}
   \left\{Z(2\lambda h)+\frac{\zeta(2+2\lambda h)^2}{b^2}\right\}.
\end{split}
\end{equation}
Its Taylor series in the two independent variables converges after
physical substitution for sufficiently large $L$.
\end{proposition}

\begin{proof}
The quotient is holomorphic at $\lambda=0$, since its numerator
vanishes there. At $h=w=0$, the right side equals
$(1-e^{-\lambda})/(\lambda\varphi(\lambda))=1$ and its $Y$
derivative is zero. The implicit function theorem therefore applies
with derivative one after moving $Y$ to the left. A finite compactness
cover and uniqueness give a common neighborhood. Substituting
$B^2=B_0^2Y$ into \eqref{jct:eq:scale} gives
\eqref{jct:eq:Y-equation} exactly. The unique large positive solution
from Theorem~\ref{jct:thm:scale} identifies the branch.
\end{proof}

\subsubsection{Exact marginal Lambert scale}

At $\eps=0$, equation \eqref{jct:eq:scale} is
\[
 \frac{bB^2}{n}=\log B+\frac32+\nu,
 \qquad \nu=\frac{\zeta(2)^2}{b^2}.
\]
Set $T=2bB^2/n$. Then
\begin{align}
 T-\log T&=L+3+2\nu-\log(2b),\label{jct:eq:Lambert-scale}\\
 T&=-W_{-1}\left(-\frac{2b}{n}e^{-3-2\nu}\right),
 &B^2&=\frac n{2b}T,\qquad q=\frac2T.
       \label{jct:eq:Lambert-explicit}
\end{align}
The lower real branch is required because $T\to\infty$
\cite{lmi:lambert}. This is a scale for the coefficient law.
The coefficient scale and the earlier inverse-function Lambert core
have different constants; the following exact matching explains them.

\begin{corollary}[All marginal central coefficients as constants]
\label{jct:cor:marginal}
Let $s=3$, $b=A=\zeta(2)$, and let $T>1$ solve
\begin{equation}\label{jct:eq:critical-T}
 T-\log T=L+5-\log(2A).
\end{equation}
Define
\begin{equation}\label{jct:eq:Dj}
 D_j=\frac1{\sqrt\pi\,j!}
   \Real\left[(\partial_z+\log2+i\pi)^j\Gamma(z)
                                  \right]_{z=j+1/2},
 \qquad D_0=1.
\end{equation}
For each fixed $N$,
\begin{equation}\label{jct:eq:marginal-T-series}
 R_b^n a_n=\sqrt{\frac A\pi}\,n^{-3/2}T^{-1/2}
       \left\{\sum_{j=0}^N\frac{D_j}{T^j}
                     +O(T^{-N-1})\right\}.
\end{equation}
If $Q_j(Y)$ denotes the Gamma-moment polynomial in
Section~\ref{lmi:section}, then exactly
\begin{equation}\label{jct:eq:matching}
 D_j=Q_j(-\log2).
\end{equation}
Thus the exact Lambert scale turns the previous logarithmic
polynomials into a sequence of constants.
\end{corollary}

\begin{proof}
At $\eps=0$ only $k=j$ survives in \eqref{jct:eq:correction}.
Since $q=2/T$, the normalized central coefficient of $T^{-j}$ is
$(-1)^j\mathcal{H}_{2j}^{[j]}(0)/(j!\mathfrak{g}(0))$.
Equation \eqref{jct:eq:central-C} gives precisely \eqref{jct:eq:Dj}.
The leading prefactor is
$\mathfrak{g}(0)/(nB)=\sqrt{A/\pi}\,n^{-3/2}T^{-1/2}$.
Finally the predecessor defines
\[
 Q_j(Y)=\frac1{\sqrt\pi\,j!}\Real\left[
             (\partial_z-Y+i\pi)^j\Gamma(z)
                                  \right]_{z=j+1/2},
\]
so substitution gives \eqref{jct:eq:matching}, with no asymptotic
comparison needed.
\end{proof}

For $C=\gamma+\log2$, the first constants are
\begin{align}
 D_1&=1-\frac C2,\notag\\
 D_2&=1-2C+\frac{3C^2}{8}-\frac{3\pi^2}{16}.
 \label{jct:eq:D-first}
\end{align}
Re-expanding \eqref{jct:eq:critical-T} in $L$ gives, in particular,
\begin{equation}\label{jct:eq:marginal-old-first}
 R_b^n a_n=\sqrt{\frac A\pi}\,n^{-3/2}L^{-1/2}
 \left\{1-\frac{\log L-\log A+\gamma+3}{2L}
              +O\left(\frac{(\log L)^2}{L^2}\right)\right\}.
\end{equation}
This agrees with the integrated marginal result, including its
constant $\gamma+3$. At higher orders the same conclusion follows
from the exact identity \eqref{jct:eq:matching} and the convergent scale
chart, rather than from a comparison of a few numerical values.

\subsubsection{Positive analytic outer functions}

The square root is not essential to the all-order logarithmic
calculus. Let $G$ vary in a family of holomorphic functions satisfying
the following uniform hypotheses for all sufficiently small real
$\eps$:
\begin{enumerate}[label=(\roman*)]
\item $G(0)=0$, and the Taylor coefficients of $G$ at zero are
nonnegative real numbers;
\item there are $\eta,M>0$ such that $G$ is holomorphic on
$|w|<\zeta(3+\eps)+\eta$ and bounded by $M$ there;
\item $G'(0)\ge\eta$.
\end{enumerate}
Shrinking $\eta$ makes it harmless to use the same symbol for the
uniform lower bound and the disk margin. Compact families with fixed
positive margins satisfy these conditions. Put
\begin{equation}\label{jct:eq:G-parameters}
 f(z)=G(\Li_{3+\eps}(z)),\quad
 c=G'(\zeta(3+\eps)),\quad
 \kappa=\frac{A^2G''(\zeta(3+\eps))}{c},\quad \mu=cA.
\end{equation}
Then $c$ is bounded above and away from zero, and $\kappa$ is
nonnegative and bounded. Define $X_n$ using \eqref{jct:eq:pgf} and set
\begin{equation}\label{jct:eq:G-scale}
 \frac{B^2}{nc}=K_\eps(\log B)+\kappa=:V,
 \qquad q=\frac{E_2(\eps)B^{-\eps}}{V},\qquad
 \xi_m=\frac{m-n\mu}{B}.
\end{equation}

\begin{theorem}[Two-derivative composition principle]
\label{jct:thm:universality}
Under the preceding uniform hypotheses, Theorem~\ref{jct:thm:all-orders}
holds with \eqref{jct:eq:G-scale} in place of
\eqref{jct:eq:scale}--\eqref{jct:eq:q-xi}, and with precisely the same
functions $\mathcal C_j$. The leading scale is
\begin{equation}\label{jct:eq:G-leading}
 B^2\sim\frac{ncL}{2}\varphi(\lambda).
\end{equation}
For the germ $v(r)=f(re^{v(r)})$, write $a_n=[r^n]v(r)$ and
$R=e^{-f(1)}$. Then $nBR^n a_n$ has the same expansion at
$\xi=(n-n\mu)/B$.

After normalization by $f(1)$ and centering, all inverse-logarithmic
orders depend on the outer function only through its first two
derivatives at $\zeta(3+\eps)$. Higher derivatives contribute to the
remainder beyond every fixed inverse-logarithmic order.
\end{theorem}

\begin{proof}
The positive Taylor coefficients and the linear term of $G$ imply
$[z^k]f(z)\ge\eta k^{-3-\eps}$. Uniform boundedness on the larger
disk supplies uniform endpoint derivatives by Cauchy's estimates.
Writing $D=\zeta(s)-\Li_s(e^{-t})$, Taylor expansion at the endpoint
gives
\[
 G(\zeta(s)-D)-G(\zeta(s))
 =-cD+\frac12G''(\zeta(s))D^2+O(|D|^3).
\]
The proof of Lemma~\ref{jct:lem:local} therefore yields
\[
 f(e^{-t})-f(1)+\mu t
  =\frac{ct^2}{2}\{K_\eps(-\operatorname{Log} t)+\kappa\}
                                      +O(|t|^{5/2}).
\]
The scale existence proof applies with $b$ replaced by $1/c$ and
$\nu$ replaced by $\kappa$. The positive-intensity proof of
Lemma~\ref{jct:lem:fourier-tail} also applies verbatim with the indicated
uniform lower bound. The entire finite-expansion proof depends only
on these two facts, so it gives the same $\mathcal C_j$ and remainder.
Lagrange inversion proves the last coefficient identity exactly as in
Proposition~\ref{jct:prop:poisson}. Terms involving higher derivatives of $G$ begin
at $O(D^3)$ and are covered by the uniform local remainder.
\end{proof}

For the square-root family, choose the disk margin smaller than
$b_-^2/2$. Its branch point is at
$w=u^2/2=\zeta(s)+b^2/2$, and $G'(0)=1/u$ is uniformly positive.
Thus the main quadratic model satisfies the hypotheses, with uniform
boundedness on a slightly smaller common margin.

\begin{center}
\begin{tabular}{@{}lll@{}}
\toprule
Outer function $G(w)$ & $c$ & $\kappa$\\
\midrule
$a w$ & $a$ & $0$\\
$u-\sqrt{u^2-2w}$ & $b^{-1}$ & $A^2/b^2$\\
$a(e^{\beta w}-1)/\beta$ & $a e^{\beta\zeta(s)}$ & $\beta A^2$\\
$-(a/\beta)\log(1-\beta w)$ & $a/(1-\beta\zeta(s))$
 & $\beta A^2/(1-\beta\zeta(s))$\\
\bottomrule
\end{tabular}
\end{center}
Here $a>0$, $\beta>0$, and in the last row
$\beta\zeta(s)<1$ with a uniform margin. The limit $\beta=0$
recovers the linear row. Compact parameter subsets satisfying the
stated margins provide concrete examples of the theorem.

The principle does not say that these models have the same full
transseries. Their coefficients at exponential scales in $L$ can
depend on higher endpoint derivatives and on other singularities.

\subsection{The first differing endpoint derivative}
\label{jct:sec:matched-jets}

The composition principle identifies all logarithmic orders but leaves
the higher derivatives of the outer function invisible at that scale.
A comparison of two models removes the common terms exactly. It gives
a hierarchy of smaller coefficient laws, without subtracting an infinite
formal expansion or choosing a summation for that expansion.

\begin{theorem}[Transport of a matched endpoint jet]
\label{jct:thm:matched-jets}
Fix an integer $r\ge3$. Let $G_1,G_2$ satisfy the uniform positivity
and analyticity hypotheses of Theorem~\ref{jct:thm:universality}, and suppose
at $z_s=\zeta(3+\eps)$ that
\[
 G_1^{(\ell)}(z_s)=G_2^{(\ell)}(z_s)\quad(1\le\ell<r),
 \qquad \delta_r=G_1^{(r)}(z_s)-G_2^{(r)}(z_s).
\]
Their endpoint values need not agree. Define $f_i=G_i\circ\Li_s$ and
the laws $\E z^{X_{i,n}}=\exp\{n(f_i(z)-f_i(1))\}$.
They have the common parameters $c,\kappa,\mu,B,q$ in
\eqref{jct:eq:G-parameters}--\eqref{jct:eq:G-scale}.
For each fixed $N\ge0$, uniformly in the compact moving-order window
and all integers $m$,
\begin{equation}\label{jct:eq:matched-jets}
\begin{split}
 B\{\mathbb{P}(X_{1,n}=m)-\mathbb{P}(X_{2,n}=m)\}
 ={}&\frac{(-1)^r A^r\delta_r n}{r!B^r}
       \sum_{j=0}^N\partial_\xi^r\mathcal C_j(\xi_m;q,\eps)\\
 &+O\!\left(\frac n{B^r}L^{-N-1}\right).
\end{split}
\end{equation}
The exact finite correction functions are
\begin{align}
 \partial_\xi^r\mathcal C_0(\xi)&=\mathcal{H}_r(\xi),\notag\\
 \partial_\xi^r\mathcal C_j(\xi;q,\eps)
 &=\frac1{j!}\sum_{k=1}^j\jctstir{j}{k}
       \left(-\frac q2\right)^k\eps^{j-k}
       \mathcal{H}_{2k+r}^{[j]}(\xi)\quad(j\ge1).
       \label{jct:eq:shifted-kernels}
\end{align}
For $v_i(\rho)=f_i(\rho e^{v_i(\rho)})$, $a_{i,n}=[\rho^n]v_i$
and $R_i=e^{-f_i(1)}$, the left side at $m=n$ is exactly
$nB(R_1^n a_{1,n}-R_2^n a_{2,n})$.
\end{theorem}

\begin{proof}
Set $a_n^*=n/B^r$, and write $\Psi_i(y)=\log\chi_{i,n}(y)$
for the explicitly centered Poisson exponent, not an unspecified branch
of the logarithm of its exponential. With $t=-iy/B$ and
$D=\zeta(s)-\Li_s(e^{-t})$, endpoint Taylor expansion gives
\begin{equation}\label{jct:eq:matched-exponent}
 \Psi_1(y)-\Psi_2(y)
 =\frac{n(-1)^r\delta_r}{r!}D^r+O(n|D|^{r+1}).
\end{equation}
The uniform local estimate proved above has $D=At+O(|t|^{3/2})$.
Consequently, for $|t|$ sufficiently small,
\[
 \frac{\Psi_1-\Psi_2}{a_n^*}
 =\frac{(-1)^r\delta_r A^r}{r!}(-iy)^r
       +O(B^{-1/2}|y|^{r+1/2}),
 \qquad |\Psi_1-\Psi_2|\le C a_n^*|y|^r.
\]

For $0\le u\le1$, the convex combination
$G_u=(1-u)G_2+uG_1$ obeys the same uniform hypotheses and has exactly
the same first two endpoint derivatives. Its centered characteristic
function is $\chi_u=\exp\{(1-u)\Psi_2+u\Psi_1\}$, and
\begin{equation}\label{jct:eq:convex-difference}
 \chi_1-\chi_2=(\Psi_1-\Psi_2)\int_0^1\chi_u\,du.
\end{equation}
The positive-intensity Fourier bound, uniformly in $u$, therefore gives
\[
 \frac{|\chi_1-\chi_2|}{a_n^*}
 \le C|y|^r e^{-c y^2}\qquad(|y|\le Bn^{-1/4}).
\]
Beyond this range, the separate bounds on $\chi_i$ give total integral
at most $2\pi B\,2e^{-c\sqrt n}/a_n^*$, smaller than every fixed
inverse power of $L$. Choose $T_L=\sqrt{D_{N,r}\log L}$ with a
sufficiently large fixed constant. The remaining tail
$T_L<|y|\le Bn^{-1/4}$ is $O(L^{-N-1})$ after division by $a_n^*$.
Thus the comparison retains its smaller scale on the whole contour;
subtracting the two unscaled remainder estimates would not suffice.

On $|y|\le T_L$, each $\chi_u$ has the same quadratic core
\eqref{jct:eq:core-characteristic}, with local error
$O(nB^{-5/2}|y|^{5/2})$, uniformly in $u$. Apply the finite exponential
and divided-exponential expansions from the all-order proof, now with
the extra integrable weight $|y|^r$. They give an $L^1$ error
$O(L^{-N-1})$. The additional errors $O(B^{-1/2})$ and
$O(nB^{-5/2})$ are smaller than every such order. The Gaussian controls
all fixed polynomial weights, while near zero the powers remain
integrable for sufficiently small $|\eps|$. Each resulting Fourier
integral contains $(-iy)^r$, so differentiation in $\xi$ shifts
$\mathcal{H}_{2k}^{[j]}$ to $\mathcal{H}_{2k+r}^{[j]}$.
The phase $e^{-i\xi_m y}$ has modulus one, proving uniformity in $m$.
Finally apply Lagrange inversion separately to each normalized model.
\end{proof}

This result isolates the first discrepancy after matching $r-1$
endpoint derivatives. Its scale is $n/B^r\asymp
n^{1-r/2}L^{-r/2}$ in the scaled probability. It does not assert a
complete transseries for a single model or a uniform expansion as $r$
or $N$ grows. The absolute assertion also permits $\delta_r=0$.

\begin{corollary}[The cubic discrepancy at the critical center]
\label{jct:cor:cubic-center}
Under the theorem's assumptions with $r=3$ and $\mu=1$,
\begin{equation}\label{jct:eq:cubic-center}
 R_1^n a_{1,n}-R_2^n a_{2,n}
 =-\frac{A^3\delta_3 q}{3B^4}
       +O(B^{-4}L^{-2}).
\end{equation}
In particular this is an asymptotic equivalence if $|\delta_3|$ is
bounded away from zero. The first nonzero central term has order
$n^{-2}L^{-3}$, although the leading shifted Gaussian vanishes there.
\end{corollary}

\begin{proof}
The exact central formula gives
$\mathcal{H}_3(0)=0$ and $\mathcal{H}_5^{[1]}(0)=-4$.
The latter follows by differentiating
$2^{(a-1)/2}/\Gamma((1-a)/2)$ at $a=5$; the derivative of reciprocal
Gamma at $-2$ is $2$. Hence
$\partial_\xi^3\mathcal C_1(0)=2q$.
Use $N=1$ in \eqref{jct:eq:matched-jets} and divide by $nB$.
Since $q\asymp L^{-1}$, the error has smaller order as asserted.
\end{proof}

\subsubsection{A positive test family}
Put $z_s=\zeta(s)$, $d_s=1+2z_s+3z_s^2$, and
\begin{equation}\label{jct:eq:positive-cubic-family}
 G_\theta(w)=\frac{w+w^2+w^3+
               \theta\{(w-z_s)^3+z_s^3\}}{A d_s},
 \qquad 0\le\theta\le\frac1{6z_s}.
\end{equation}
All three Taylor coefficients are strictly positive and $G_\theta(0)=0$.
At the endpoint, $c=1/A$, $\mu=1$ and
$\kappa=A^2(2+6z_s)/d_s$ are independent of $\theta$, whereas
$G_\theta'''-G_0'''=6\theta/(A d_s)$.
Thus increasing $\theta$ makes the normalized diagonal coefficient
eventually smaller, with the explicit leading difference
\eqref{jct:eq:cubic-center}. Each coefficient must use its own $R_\theta$;
the endpoint values differ by $\theta z_s^3/(A d_s)$.

Finite positive coefficient recurrences provide a direct experimental
test without using Fourier asymptotics to compute the coefficients.
If $f(z)=\sum_{k\ge1}f_kz^k$, compute
\[
 p_0=e^{-nf(1)},\qquad
 p_m=\frac n m\sum_{k=1}^m k f_k p_{m-k},
 \qquad R^n a_n=\frac{p_n}{n}.
\]
The preserved diagnostics use these recurrences for both models and
compare their difference with the shifted-kernel formula. Such finite
floating-point comparisons guide and check the identities; the
uniform remainder is supplied by the contour argument above.


\subsection{A growing positive window and its Lambert transition}
\label{jct:sec:positive-wing}

The absolute error in Theorem~\ref{jct:thm:all-orders} is uniform in the
lattice location.  Its usefulness is therefore not confined to bounded
standardized displacements.  We now use its explicit kernels to obtain
a relative approximation in growing positive windows.

Throughout this section, put $L=\log n$ and retain the scale $B$, the
parameters $q,\eps$, and the centered coordinate
$\xi=(m-nA/b)/B$ of Theorem~\ref{jct:thm:all-orders}.  Write
\[
 \mathcal P_n(\xi)=B\mathbb P(X_n=m),\qquad
 g(\xi)=\mathfrak{g}(\xi)=\frac{\exp(-\xi^2/2)}{\sqrt{2\pi}}.
\]
Only lattice coordinates are meant when $\mathcal P_n$ occurs.  The
uniform parameter hypotheses of that theorem give
\begin{equation}\label{jct:wing:small-parameters}
 \frac{c}{L}\le q\le\frac{C}{L},\qquad
 |\eps|\le\frac{C}{L},
\end{equation}
with positive constants independent of the parameters.  In particular,
the constants in the statements below are uniform for the scaled order
in a fixed compact real set and for $b$ in a compact subinterval of
$(0,\infty)$.

\subsubsection{Right tails of the explicit kernels}

We use the notation
\[
 \mathcal H_{2k}^{[j]}(x)
 =\left.\frac{\partial^j}{\partial\alpha^j}
             \mathcal H_\alpha(x)\right|_{\alpha=2k},\qquad
 \mathcal H_\alpha(x)=\frac1{2\pi}
       \int_{\mathbb R}e^{-iyx-y^2/2}(-iy)^\alpha\,dy,
\]
where $\operatorname{Re}\alpha>-1$ and the principal powers are used.

\begin{samepage}
\begin{lemma}[Kernel tails and order derivatives]\label{jct:lem:kernel-wing}
For fixed integers $k\ge1$ and $j\ge1$, as $x\to+\infty$,
\begin{equation}\label{jct:wing:derivative-bound}
 \mathcal H_{2k}^{[j]}(x)
 =O\bigl(x^{-2k-1}(1+\log x)^{j-1}\bigr).
\end{equation}
More precisely, if $a(\alpha)=1/\Gamma(-\alpha)$, then
\begin{equation}\label{jct:wing:leading-polynomial}
 \mathcal H_{2k}^{[j]}(x)
 =x^{-2k-1}\sum_{r=1}^j\binom jr
        a^{(r)}(2k)(-\log x)^{j-r}
       +O\bigl(x^{-2k-11/4}\bigr).
\end{equation}
In particular,
\begin{equation}\label{jct:wing:first-kernel-tail}
 \mathcal H_2^{[1]}(x)
 =-2x^{-3}\left\{1+6x^{-2}+O(x^{-4})\right\}.
\end{equation}
The constants can depend on $k$ and $j$, which are fixed here.
\end{lemma}
\end{samepage}

\begin{proof}
We give a convolution proof so that differentiation with respect to
the order does not rest on an unstated uniformity of a special-function
asymptotic formula.  Fix $k$, set $d=2k+1$, and let
$|\alpha-2k|\le1/4$.  Then
\begin{equation}\label{jct:wing:convolution}
 \mathcal H_\alpha(x)
 =\frac1{\Gamma(d-\alpha)}\int_0^\infty
               t^{d-\alpha-1}g^{(d)}(x-t)\,dt.
\end{equation}
To verify the identity, first insert $e^{-\eta t}$ with $\eta>0$
in the integral before differentiating $d$ times in $x$.  The Fourier
multiplier of the resulting convolution is
$(\eta-iy)^{\alpha-d}$.  Differentiation multiplies it by
$(-iy)^d$.  As $\eta\downarrow0$, this tends to $(-iy)^\alpha$.
It is dominated, up to a constant uniform on the indicated order
disc, by $|y|^{\operatorname{Re}\alpha}e^{-y^2/2}$: here
$\operatorname{Re}(\alpha-d)<0$ and $|\eta-iy|\ge|y|$.
This is integrable.  Dominated convergence on the physical side is
also valid, because the power at $t=0$ has real part greater than
$-1$ and every derivative of $g$ has Gaussian decay at infinity.
This proves \eqref{jct:wing:convolution}, jointly holomorphically in
$\alpha$ on a neighborhood of the closed disc.

Put $p=d-\alpha-1$ and change variables to $u=t-x$.  On
$|u|\le x/2$, expand $(x+u)^p$ in powers of $u/x$ through degree
$d+2M+1$, for any fixed $M\ge0$.  The Taylor remainder is bounded by
\[
 C_{k,M}x^{\operatorname{Re}p-d-2M-2}|u|^{d+2M+2}
\]
uniformly on the order disc.  Integration against
$|g^{(d)}(-u)|$ is harmless.  Replacing the central integrals by
integrals over the whole real line changes the answer by a quantity
bounded by $e^{-x^2/16}$ times a fixed power of $x$.
The part of the original integral where $|u|>x/2$ has the same
type of bound.  For the left part, the integrable factor
$(x+u)^{\operatorname{Re}p}$ near $u=-x$ causes no difficulty;
for the right part, polynomial growth is dominated by the Gaussian.

Integration by parts gives
\[
 \int_{\mathbb R}u^r g^{(d)}(-u)\,du=0
 \quad(r<d\text{ or }r-d\text{ odd}),
\]
and, for $r=d+2\ell$,
\[
 \int_{\mathbb R}u^{d+2\ell}g^{(d)}(-u)\,du
 =\frac{(d+2\ell)!}{2^\ell\ell!}.
\]
It follows that, uniformly on the order disc,
\begin{equation}\label{jct:wing:uniform-tail-expansion}
 \mathcal H_\alpha(x)
 =\sum_{\ell=0}^{M}
       \frac{x^{-\alpha-2\ell-1}}
            {2^\ell\ell!\,\Gamma(-\alpha-2\ell)}
       +O\bigl(x^{-\operatorname{Re}\alpha-2M-3}\bigr).
\end{equation}
The exponentially small terms just described can be retained
separately; they are already absorbed by this remainder.

The remainder in \eqref{jct:wing:uniform-tail-expansion} is
holomorphic in $\alpha$.  Cauchy's formula on the circle of radius
$1/4$ about $2k$ bounds any fixed order derivative of that remainder
by $O(x^{-2k-2M-11/4})$.  With $M=0$, differentiating the leading
term gives \eqref{jct:wing:leading-polynomial}: the $r=0$ term vanishes
because $a(2k)=0$.  This proves \eqref{jct:wing:derivative-bound}.
For \eqref{jct:wing:first-kernel-tail}, use $M=2$ and $k=j=1$.
The reciprocal Gamma function has a simple zero at each
nonpositive integer, and
\[
 \left.\frac{d}{d\alpha}\frac1{\Gamma(-\alpha-2\ell)}
                         \right|_{\alpha=2}
 =-(2+2\ell)!.
\]
The first two terms are therefore $-2x^{-3}$ and $-12x^{-5}$;
the remaining displayed term and the differentiated remainder are
$O(x^{-7})$, as asserted.
\end{proof}

The same tail structure follows from the classical connection formula
for parabolic cylinder functions and their large-variable expansions \cite{jct:cylinder}:
$\mathcal H_\alpha(x)=e^{-x^2/4}D_\alpha(-x)/\sqrt{2\pi}$,
\[
 D_\alpha(-x)=\cos(\pi\alpha)D_\alpha(x)
       +\frac{\pi}{\Gamma(-\alpha)}V(-\alpha-1/2,x).
\]
Here the first term contributes Gaussian decay and the second
contributes the power tail.  The convolution proof above supplies
the parameter uniformity and the order derivatives directly.

\subsubsection{A relative Gaussian and power-tail approximation}

\begin{theorem}[Positive windows of polynomial size in $\log n$]
\label{jct:thm:positive-windows}
Fix $M>0$, and let $r_L\to\infty$ with $r_L\le L^M$.
Uniformly over lattice points and admissible parameters for which
$r_L\le\xi\le L^M$,
\begin{equation}\label{jct:wing:relative-mixture}
 \mathcal P_n(\xi)
 =\left\{g(\xi)+q\xi^{-3}\right\}\{1+o(1)\}.
\end{equation}
Consequently, for each fixed $\delta\in(0,2)$, whenever the stated
intervals are nonempty,
\begin{align}
 \mathcal P_n(\xi)&\sim g(\xi),
 &r_L\le\xi\le\sqrt{(2-\delta)\log L},
                       \label{jct:wing:gaussian-window}\\
 \mathcal P_n(\xi)&\sim q\xi^{-3},
 &\sqrt{(2+\delta)\log L}\le\xi\le L^M.
                       \label{jct:wing:power-window}
\end{align}
No relative conclusion about negative growing displacements is
included in this theorem.
\end{theorem}

\begin{proof}
Choose an integer $N>3M$ in Theorem~\ref{jct:thm:all-orders}.
Its term of total degree one is
$-q\mathcal H_2^{[1]}(\xi)/2$.  By
\eqref{jct:wing:first-kernel-tail}, this equals
$q\xi^{-3}\{1+O(\xi^{-2})\}$.
For every term with total degree $j\ge2$, divide its absolute
value by $q\xi^{-3}$.  Equations
\eqref{jct:wing:small-parameters} and \eqref{jct:wing:derivative-bound}
bound the quotient by a constant times
\[
 q^{k-1}|\eps|^{j-k}
       \xi^{-2k+2}(1+\log\xi)^{j-1}
 \le C_N\left(\frac{1+\log\xi}{L}\right)^{j-1}.
\]
Only finitely many terms occur.  Their sum is $o(1)$ uniformly
because $\log\xi\le M\log L$.
The absolute remainder divided by $q\xi^{-3}$ is
\[
 O\bigl(\xi^3L^{-N}\bigr)=O\bigl(L^{3M-N}\bigr)=o(1).
\]
Finally $\xi^{-2}\le r_L^{-2}\to0$.
We have thus proved the stronger additive statement
$\mathcal P_n(\xi)=g(\xi)+q\xi^{-3}\{1+o(1)\}$,
which implies \eqref{jct:wing:relative-mixture} since both main terms
are positive.

For \eqref{jct:wing:gaussian-window}, the ratio of the power term to
the Gaussian is at most a constant times
$L^{-\delta/2}r_L^{-3}$.  For \eqref{jct:wing:power-window}, the
ratio of the Gaussian to the power term is at most a constant
times $L^{-\delta/2}(\log L)^{3/2}$: the function
$x^3e^{-x^2/2}$ decreases for $x>\sqrt3$.
Both bounds tend to zero, proving the two consequences.
\end{proof}

In the coefficient specialization of Theorem~\ref{jct:thm:all-orders},
$m=n$, and hence the same conclusions apply directly to
$nBR_b^n a_n$.  Coordinates $\xi\le L^M$ force $b\to A$
because $1-A/b=\xi B/n=O(L^{M+1/2}/\sqrt n)$.
Thus these are growing positive windows adjacent to the feedback
threshold.  They are distinct from a fixed supercritical parameter
regime, whose standardized displacement is much larger.

\begin{corollary}[Relative matching with the supercritical coefficient law]
\label{jct:cor:supercritical-matching}
In the diagonal specialization, uniformly for each fixed $M>0$ and
$\delta>0$ in the nonempty range
\[
 \sqrt{(2+\delta)\log L}\le\xi\le L^M,
\]
one has
\begin{equation}\label{jct:eq:supercritical-matching}
 R_b^n a_n\sim
 \frac1{b\{n(1-A/b)\}^{3+\eps}}
 =\frac{b^{s-1}}{(b-A)^s}\,n^{-s}.
\end{equation}
\end{corollary}
\begin{proof}
On this window $E_2(\eps)\to1$ and
$\xi^{-\eps}\to1$ uniformly, since
$|\eps\log\xi|=O(\log L/L)$. Thus the power term in
\eqref{jct:wing:power-window} is relatively equivalent to
$q\xi^{-3-\eps}/E_2(\eps)$. The exact identities
\[
 \frac{q}{E_2(\eps)}=\frac{n}{bB^{2+\eps}},\qquad
 B\xi=n(1-A/b)
\]
give \eqref{jct:eq:supercritical-matching} after division by $nB$.
\end{proof}

There is an exact explanation for the amplitude cancellation.
The first divided-difference correction is
$-q(\mathcal{H}_{2+\eps}-\mathcal{H}_2)/(2\eps)$.
Its leading algebraic tail is
$q\xi^{-3-\eps}/E_2(\eps)$ because
$1/\Gamma(-2-\eps)=-2\eps/E_2(\eps)$.
The polynomial-log window makes this term relatively equivalent to
the simpler $q\xi^{-3}$ while preserving the original coefficient
amplitude in \eqref{jct:eq:supercritical-matching}.

\subsubsection{An exact Lambert location and a two-term profile}

\begin{theorem}[Lambert location of the Gaussian--power transition]
\label{jct:thm:wing-lambert}
For sufficiently small $q>0$, let $x_*(q)>\sqrt3$ be the larger
positive solution of $g(x_*)=q x_*^{-3}$.  It has the exact formula
\begin{equation}\label{jct:wing:lambert-center}
 x_*(q)^2
 =-3W_{-1}\left(-\frac{(2\pi)^{1/3}}3q^{2/3}\right).
\end{equation}
If $Q=\log(1/q)$, then
\begin{equation}\label{jct:wing:expanded-center}
 x_*(q)^2
 =2Q+3\log(2Q)-\log(2\pi)
       +O\left(\frac{\log Q}{Q}\right).
\end{equation}
For the parameters in \eqref{jct:wing:small-parameters}, uniformly
for $\tau$ in compact real intervals and positive lattice coordinates
$\xi_n>0$ satisfying
\begin{equation}\label{jct:wing:transition-coordinate}
 \xi_n^2=x_*(q)^2+2\tau+o(1),
\end{equation}
where the final error is uniform, one has
\begin{equation}\label{jct:wing:transition-profile}
 \frac{\mathcal P_n(\xi_n)}{qx_*(q)^{-3}}
       \longrightarrow 1+e^{-\tau}.
\end{equation}
In particular $x_*(q)\sim\sqrt{2\log\log n}$.
\end{theorem}

\begin{proof}
The function $x^3g(x)$ increases up to $\sqrt3$ and decreases
thereafter, so the stated larger root exists and is unique for
small $q$.  The defining equation is equivalent to
\[
 x_*^2e^{-x_*^2/3}=(2\pi)^{1/3}q^{2/3},
\]
which gives \eqref{jct:wing:lambert-center} with the branch $W_{-1}$ \cite{lmi:lambert}.
Alternatively, taking logarithms gives the exact scalar equation
\begin{equation}\label{jct:wing:center-equation}
 y=2Q+3\log y-\log(2\pi),\qquad y=x_*^2.
\end{equation}
It first implies $y\sim2Q$.  Substituting
$y_0=2Q+3\log(2Q)-\log(2\pi)$ into its two sides leaves a
residual $O((\log Q)/Q)$, while the derivative of
$y-3\log y$ is $1-3/y$, bounded away from zero at the larger
root.  The mean value theorem proves
\eqref{jct:wing:expanded-center}.

By \eqref{jct:wing:small-parameters}, $Q=\log L+O(1)$.
Therefore \eqref{jct:wing:transition-coordinate} lies in the
window of Theorem~\ref{jct:thm:positive-windows} for any fixed $M>0$
and a suitable $r_L\to\infty$.  Moreover,
\[
 \frac{g(\xi_n)}{qx_*(q)^{-3}}
 =\frac{g(\xi_n)}{g(x_*(q))}
 =e^{-\tau+o(1)},\qquad
 \frac{q\xi_n^{-3}}{qx_*(q)^{-3}}
 =\left(\frac{x_*(q)}{\xi_n}\right)^3\longrightarrow1.
\]
The relative mixture theorem now gives
\eqref{jct:wing:transition-profile}.
\end{proof}

The lattice restriction is immaterial to the transition coordinate:
rounding $nA/b+B\sqrt{x_*(q)^2+2\tau}$ to the nearest
integer changes $\xi_n^2$ by $O(x_*(q)/B)=o(1)$.
For diagonal coefficients, the feedback parameter can instead be
varied on the locally unique branch near $b=A$, choosing the compact
feedback interval to contain $\zeta(2)$ in its interior.  Here is an explicit
verification of the latter assertion.  Hold $n,s$ fixed, put
$V=bB^2/n$, and differentiate the defining scale equation.  It gives
\[
 \beta:=\partial_b\log B
 =-\frac{b^{-1}+2A^2/(b^3V)}{2-q}=O(1),\qquad
 \partial_b\log q=-b^{-1}-(2+\varepsilon)\beta=O(1).
\]
For $\xi(b)=n(1-A/b)/B$, consequently,
\[
 \xi'(b)=\frac{nA}{b^2B}-\xi(b)\beta
                  \asymp\sqrt{n/L}>0
\]
through any window $\xi=O(L^M)$ near $b=A$.  The derivative of
$x_*(q(b))$ is only $O(1/x_*)$, by
\eqref{jct:wing:center-equation}.  Starting from $\xi(A)=0$,
continuity and these derivative bounds yield a unique local
feedback choice for each bounded $\tau$ in
\eqref{jct:wing:transition-coordinate}.

\subsubsection{Relation to established probability results}

The Gaussian and power-tail competition is an established phenomenon.
Lam, Blanchet, Burch, and Bazant obtain nonclassical corrections to
Gaussian limits for regularly varying densities, express corrections
using Dawson and parabolic cylinder functions, and explain their
connection with heavy-tail large deviations~\cite{jct:lam}.
Berger, Birkner, and Yuan treat large and local large deviations,
including the infinite-variance normal domain with tail index $-2$,
and identify the transition between collective Gaussian behavior and
a single exceptionally large summand~\cite{jct:berger}.

Accordingly, neither the existence of such a competition nor the use
of parabolic cylinder functions is claimed as new here.  The result
established for the present model is the uniform extension through the
varying order $s_n=3+O(1/\log n)$ and varying nonlinear feedback,
with the explicit paired normalization and the finite all-order
Stirling formulas of Theorem~\ref{jct:thm:all-orders}.  The growing windows
above follow from its uniform absolute remainder, and the exact
Lambert coordinate \eqref{jct:wing:lambert-center} describes their
Gaussian--power transition in that normalization.  The fixed-law
results just cited do not, without additional uniform estimates,
supply this triangular-parameter assertion.

\section{Reciprocal orders and removable inverse endpoints}
\label{jct:sec:endpoints}

The coefficient transition concerns order three. A different exceptional
set appears when directly inverting a polylogarithm at its finite
endpoint: the reciprocal orders $s=1+1/m$. We develop their convergent
moving-order chart after recalling the local ramification classification.
That classification and the cut-modulus obstruction also occur in the
earlier reversion continuation \cite{llc:reversion}; the joint report
\cite{jct:report} gives the proof route used here and the uniform chart.
The direct-inverse predecessor \cite{jct:inverse} states its global
endpoint-circle theorem only for integer orders. The obstruction below
limits its proposed real-order extension and does not retract that theorem.

The nonintegral charts in the predecessor imply an exceptional phenomenon
that is worth separating from the integer-order resonance theory.
Write $s=1+\alpha>1$ and
\[
 D_\alpha(t)=\zeta(1+\alpha)-\operatorname{Li}_{1+\alpha}(e^{-t}).
\]
For nonintegral $\alpha$, the classical convergent expansion \cite{rw:DLMFPolylog} gives
\[
 D_\alpha(t)=b(\alpha)t^\alpha+tA_\alpha(t),\qquad
 b(\alpha)=-\Gamma(-\alpha),\qquad
 A_\alpha(t)=\sum_{k\ge1}\frac{(-1)^{k+1}\zeta(1+\alpha-k)}{k!}t^{k-1}.
\]
Initially $0<|t|<2\pi$, $-\pi<\arg t<\pi$, and
$t^\alpha=\exp(\alpha\log t)$ uses the principal logarithm.
Further continuation is on the logarithmic covering.
The function $A_\alpha(t)$ is jointly holomorphic for $t$ near zero and
$\alpha$ in any sufficiently small neighborhood avoiding the nonnegative
integers.  Let $t_\alpha(y)$ denote the inverse which is positive for
positive small $y$.

\begin{theorem}[Minimal local ramification]\label{jct:thm:ramification}
For real $s>1$ the minimal degree of a finite ramification $y=w^N$
making the endpoint inverse $t_{s-1}(y)$ holomorphic is as follows.
\begin{enumerate}
\item If $0<\alpha=s-1<1$ and $\alpha=p/q$ is in lowest terms,
the minimal degree is $p$.
\item If $\alpha>1$ is nonintegral and $\alpha=p/q$ is in lowest
terms, the minimal degree is $q$.
\item If $\alpha$ is irrational, or if $s$ is a positive integer
at least two, no finite ramification suffices.
\end{enumerate}
In particular, the endpoint inverse is holomorphic without ramification
if and only if $s=1+1/m$ for an integer $m\ge2$.
The same classification holds for the endpoint germ
$\operatorname{Li}_{s}^{\langle-1\rangle}(\zeta(s)-y)=e^{-t_{s-1}(y)}$.
\end{theorem}

\begin{proof}
Suppose first that $0<\alpha=p/q<1$.  The substitution $t=\xi^q$
makes the forward germ holomorphic:
\[
 D_\alpha(\xi^q)=b\xi^p+\xi^q A_\alpha(\xi^q).
\]
The factor multiplying $\xi^p$ is nonzero at zero, so an analytic
change of coordinate reduces this germ to $\eta^p$.  Its local inverse
is therefore holomorphic in $y^{1/p}$, and its $t$ coordinate starts
with $b^{-q/p}y^{q/p}$.  Since $p$ and $q$ are coprime, this leading
term forces every possible ramification degree to be divisible by $p$.
Thus the minimal degree is exactly $p$.

If $\alpha=p/q>1$ is nonintegral, the same substitution gives a
forward germ of order $q$, because $A_\alpha(0)=\zeta(\alpha)>0$.
It follows that the inverse $t$ is holomorphic in $y^{1/q}$.  To see
minimality, let $T(y)$ be the analytic inverse of $tA_\alpha(t)$.
Expansion of the inverse equation gives
\[
 t_\alpha(y)=T(y)-b\zeta(\alpha)^{-\alpha-1}y^\alpha+o(y^\alpha),
\]
after retaining the analytic terms of $T$ through the relevant degree.
The displayed fractional coefficient is nonzero, so $q$ divides every
possible ramification degree.

For irrational $\alpha<1$, the leading exponent $1/\alpha$ is
irrational.  For irrational $\alpha>1$, the preceding first
nonanalytic exponent is $\alpha$.  Neither can occur in a convergent
Puiseux germ.  At $s=2$,
$D_1(t)=t(1-\log t)+O(t^2)$, whence
$t_1(y)\sim y/\log(1/y)$, which is not Puiseux.
At integer $s=m+1\ge3$, the singular term of the forward deficit is
$(-1)^m t^m\log t/m!$, and its linear coefficient is $a=\zeta(m)>0$.
Therefore the inverse has a nonzero logarithmic term
$(-1)^{m+1}y^m\log y/(m!a^{m+1})$ in addition to its analytic part.
Finite ramification multiplies the logarithm by a nonzero integer
and does not remove it.
Finally $t\mapsto e^{-t}$ is biholomorphic near zero, so it preserves
these ramification degrees.
\end{proof}

\subsection{Explicit coefficients at the removable endpoints}

Fix an integer $m\ge2$, put $\alpha=1/m$ and $b_m=-\Gamma(-1/m)>0$,
and define
\[
 G_m(\xi)=b_m+\sum_{k\ge1}
 \frac{(-1)^{k+1}\zeta(1+1/m-k)}{k!}\xi^{mk-1}.
\]
Then $y=\xi G_m(\xi)$, $t=\xi^m$, and Lagrange inversion gives
the exact finite extraction formulas
\[
 [y^n]t_{1/m}(y)=\frac mn[\xi^{n-m}]G_m(\xi)^{-n},
\]
\[
 [y^n]e^{-t_{1/m}(y)}
 =-\frac mn[\xi^{n-m}]e^{-\xi^m}G_m(\xi)^{-n}\qquad(n\ge1).
\]
Coefficients with a negative extraction exponent vanish.  In particular
\[
 t_{1/m}(y)=b_m^{-m}y^m
 -m\zeta(1/m)b_m^{-2m}y^{2m-1}+O(y^{3m-2}).
\]
This last remainder follows because the next possible degree in
$t(y)$ is $3m-2$ (the term with $k=2$ first contributes at $3m-1$).
For the ordinary inverse endpoint germ, every derivative of order
$1,\ldots,m-1$ vanishes, and its $m$th derivative with respect to
$w=\zeta(s)-y$ equals
$(-1)^{m+1}m!b_m^{-m}$.

For example, at $s=3/2$ one has $b_2=2\sqrt\pi$ and
\[
 t_{1/2}(y)=\frac{y^2}{4\pi}
 -\frac{\zeta(1/2)}{8\pi^2}y^3
 +\frac{5\zeta(1/2)^2}{64\pi^3}y^4+O(y^5).
\]
Thus the square-root singularity of the forward polylogarithm has
become a quadratic stationary point of its holomorphic endpoint inverse.

\subsection{A uniform chart crossing a removable endpoint}

The apparent disappearance of branching is not uniform under order
perturbation.  Nevertheless a short exact equation gives a convergent
chart that is uniform at the reciprocal order.

\begin{samepage}
\begin{theorem}[Reciprocal-order crossover chart]\label{jct:thm:reciprocal-chart}
Let $m\ge2$ and let $K$ be a compact real interval.  Put
\[
 \alpha=\frac1m+\epsilon,\quad T=\log(1/y),\quad
 \tau=\epsilon T,\quad h=T^{-1},\quad r=y^{m-1},\quad q=y^m.
\]
There is a function $U(h,r,q;\tau)$ holomorphic in its first three
variables near zero, with a complex neighborhood of $K$ for the last
variable and neighborhoods uniform on $K$, such that the physical
inverse satisfies
\[
 t_\alpha(y)=y^m U(T^{-1},y^{m-1},y^m;\epsilon T).
\]
Its value at the origin is
\[
 U(0,0,0;\tau)=b_m^{-m}e^{m^2\tau}.
\]
The Taylor series in the three independent small variables converges
after physical substitution for sufficiently small positive $y$.
\end{theorem}
\end{samepage}

\begin{proof}
Substitute $t=y^m U$ into $D_\alpha(t)=y$ and divide by $y$.
The exact equation is
\[
 b(1/m+h\tau)e^{-m\tau}U^{1/m+h\tau}
 +rU A_{1/m+h\tau}(qU)=1.
\]
For $h=r=q=0$ its indicated solution is
$U_0(\tau)=b_m^{-m}e^{m^2\tau}$.
Choose the logarithm
$\log U_0(\tau)=-m\log b_m+m^2\tau$ and write $U=U_0V$.
In this coordinate the equation is holomorphic near $V=1$, and its
$V$ derivative at $h=r=q=0$, $V=1$ is $1/m$.
The holomorphic implicit function theorem and a finite compactness
cover of $K$ give the claimed uniform chart.  Local uniqueness glues
the charts.  The physical positive branch is selected by continuity
from positive $y$ and the strictly increasing real forward deficit.
\end{proof}

For each fixed $j\ge1$, differentiation in $\alpha$ at fixed deficit
$y$ gives
\[
 \left.\partial_\alpha^j t_\alpha(y)\right|_{\alpha=1/m}
 =b_m^{-m}(-m^2)^j y^m(\log y)^j
 +O\bigl(y^m(1+|\log y|)^{j-1}\bigr).
\]
One may prove this either from the uniform chart or by differentiating
$b(\alpha)^{-1/\alpha}y^{1/\alpha}$ and the jointly holomorphic
two-generator remainder.  Derivatives of the latter acquire at most
fixed powers of $\log y$ and retain a positive additional power of
$y$.  Thus order differentiation exposes an infinite logarithmic
monodromy which is absent from the base germ.
Here the independent variable is the deficit $y$; differentiation at
fixed $w$ also differentiates $\zeta(s)-w$ and is a different operation.

\subsection{Why the integer endpoint-circle proof cannot extend verbatim}

For $0<\alpha=s-1<1/2$ and $u\downarrow0$, the upper-bank value is
\[
 \operatorname{Li}_{1+\alpha}(e^u+i0)
 =\zeta(1+\alpha)+\Gamma(-\alpha)e^{-i\pi\alpha}u^\alpha+O(u).
\]
Since $\Gamma(-\alpha)<0$ and $\cos(\pi\alpha)>0$, this yields
\[
 |\operatorname{Li}_{s}(e^u+i0)|
 =\zeta(s)+\Gamma(1-s)\cos(\pi(s-1))u^{s-1}
 +o(u^{s-1})<\zeta(s)
\]
for every sufficiently small positive $u$.
Consequently the cut-modulus assertion in the integer-order theorem
fails for every $1<s<3/2$.  The predecessor states that assertion only
for integers, so this is a negative answer to a proposed extension,
not a counterexample to a printed theorem.

The cut entering the endpoint circle does not itself determine the
Taylor radius of an analytically continued inverse: the cut is a
boundary of the chosen forward sheet, not necessarily an inverse
singularity.  Likewise endpoint removability does not rule out another
singularity of the same modulus.  Determining the global radius at the
reciprocal orders remains a separate problem.

\subsection{Evidence and remaining analytic questions}

The delivered joint report's three programs have a preserved isolated
replay: exact original-equation versus Lagrange coefficients through
degree eight, six finite Stirling identities, 36 central Gamma identities,
five marginal coefficient identities, nine noninteger Fourier comparisons,
36 coefficient cases through $n=4096$, nine endpoint diagnostics and six
cut-bank diagnostics. The supplied article, PDF, code and outputs remain
unchanged. The separate experiment for Theorem~\ref{jct:thm:matched-jets}
uses the positive cubic family \eqref{jct:eq:positive-cubic-family}; its
coefficient recurrence and order-kernel evaluation are independent routes.
All floating-point results are non-interval diagnostics. The infinite
claims rest on the written analytic proofs, and no proof-assistant
formalization is asserted.

The joint theorem settles compact real $(s-3)\log n/2$ windows at every
fixed logarithmic order. The positive relative theorem reaches every
fixed power of $\log n$ in the standardized coordinate. Growing negative
relative windows, unbounded moving-order parameters, and macroscopic
positive-tail matching remain open. The matched-derivative theorem
identifies the first discrepancy between suitable models; a complete
single-model expansion of all smaller sectors is a separate question.
Likewise the convergent reciprocal-order chart is local and does not
determine the global Taylor radius after endpoint removability.
